\documentclass[10pt]{article}
\usepackage[letterpaper,margin=0.62in,top=0.55in,bottom=0.6in]{geometry}
\usepackage{amsmath,amssymb,amsthm}
\usepackage[T1]{fontenc}
\usepackage{lmodern}
\usepackage{microtype}
\usepackage{titlesec}
\usepackage[hidelinks]{hyperref}
\usepackage{enumitem}

\titleformat{\section}[runin]{\normalfont\bfseries}{\thesection.}{0.4em}{}[.]
\titlespacing*{\section}{0pt}{0.55ex}{0.6em}
\setlength{\parindent}{0pt}
\setlength{\parskip}{0.45ex}
\pagestyle{empty}
\newtheorem*{thmA}{Theorem 1}
\newtheorem*{thmB}{Theorem 2}
\newcommand{\ch}[2]{\{#1,#2\}}

\begin{document}

\begin{center}
{\large\bfseries Six chords are exactly enough for $41\le n\le 67$:\\[1pt]
minimal pancyclic graphs found by AI agents (Erd\H{o}s Problem \#1016)}\\[4pt]
Joshua Pinckard, ToolsEnabled \quad\texttt{josh@toolsenabled.ai}\quad 23 September 2026\\[1pt]
{\small Working draft. The research was carried out by LLM agents (Section 4). Code, proofs and data: \url{https://github.com/JoshuaPinckard/erdos-1016-pancyclic}}
\end{center}
\vspace{-0.6ex}

\textbf{Abstract.} A graph on $n$ vertices is \emph{pancyclic} if it has a cycle of every length $3,\dots,n$. Let $m(n)=n+h(n)$ be the fewest edges of a pancyclic graph on $n$ vertices. We show $h(n)=6$ for every $41\le n\le 67$, so $m(n)=n+6$ there, and that the largest $n$ with a 6-chord pancyclic graph is $67$ or lies in $\{72,\dots,79,81,83,84,85\}$. The upper bounds are checked in Lean~4. The lower bounds are exhaustive computer searches. AI agents did the work: literature search, algorithms, GPU kernels, Lean proofs and adversarial review.

\section{Problem}
A pancyclic graph contains a Hamilton cycle, so $h(n)$ is the least number of chords whose addition to the cycle $C_n$ gives a pancyclic graph. Bondy~[1] stated $\log_2(n-1)-1\le h(n)\le \log_2 n+\log_* n+O(1)$, and a proof of the upper bound is attributed to~[6, Ch.~4.5]. Erd\H{o}s asked whether $h(n)-\log_2 n\to\infty$~[2]. Griffin~[3] determined $m(n)$ for $n\le 37$. White and Claude~[4] gave $h(n)=5$ for $34\le n\le 40$, and Robinfxa~[5] proved $m(41)=47$ with a Lean-checked exclusion. Let $t_k=\max\{n: h(n)\le k\}$. Griffin's table gives $t_1,\dots,t_4=5,8,14,24$.

\section{Results}
\begin{thmA}
$h(n)=6$ for every $41\le n\le 67$. Hence $m(n)=n+6$ on this range, and $t_5=40$.
\end{thmA}
\emph{Upper bound.} For $41\le n\le 67$, let $F_n$ be $C_n$ (vertices $0,\dots,n-1$) plus the chords
$\ch{0}{2}$, $\ch{0}{n{-}7}$, $\ch{1}{13}$, $\ch{3}{n{-}6}$, $\ch{4}{31}$, $\ch{n{-}8}{n{-}5}$.
Its cycle lengths are exactly $[3,32]\cup[n-34,n]$, which is all of $[3,n]$ precisely when $n\le 67$. $F_{68}$ misses only length $33$. The Lean~4/Mathlib theorem \texttt{family\_pancyclic\_41\_67} checks every case in the kernel, using only the axioms \texttt{propext}, \texttt{Classical.choice} and \texttt{Quot.sound}.
\emph{Lower bound.} No $C_n$ plus 5 chords is pancyclic for any $n\ge 41$. An exhaustive shape search (Section 3) refutes every $n$ from $41$ to $58$, and no 5-chord shape can be pancyclic beyond $n=58$. At $n=41$ a complete GPU enumeration of all 5-chord sets agrees. This confirms [5] by a different method.

\begin{thmB}
No 6-chord pancyclic graph exists for $n\in\{68,69,70,71,80,82\}$ or for any $n\ge 86$. Hence $t_6=67$ or $t_6\in\{72,\dots,79,81,83,84,85\}$.
\end{thmB}
A CPU branch-and-bound refutes every $n$ from $88$ to $111$, and $111$ is the largest $n$ any 6-chord shape can reach. GPU exhaustion refutes $68$--$71$, $80$, $82$ and $86$--$89$, so $88$ and $89$ are refuted twice, by two different programs. The runs stopped on 20 September. Levels $84$ and $85$ each lack one tier. A blind control run over all of level $67$ recovered $F_{67}$ at its expected shape and also found a second 67-vertex witness, not isomorphic to $F_{67}$: chords $\ch{0}{27}$, $\ch{18}{30}$, $\ch{28}{37}$, $\ch{29}{37}$, $\ch{31}{38}$, $\ch{32}{39}$.

\section{Method}
$C_n$ plus $k$ chords is a subdivision of one of finitely many multigraphs, called \emph{shapes}: $1{,}236$ for $k=5$ and $21{,}878$ for $k=6$. Each cycle is a fixed set of arcs, so its length is a sum of arc lengths, and pancyclicity becomes a covering problem in the arc lengths. A shape with $F$ distinct cycle forms can be pancyclic only if $n\le F+2$, which bounds every search. The CPU solver decides each pair of shape and $n$ by branch and bound. The GPU pipeline enumerates arc-length compositions. It prunes them with a pairwise interval-Hall test, which is necessary for pancyclicity and monotone in the arc bounds, and which keeps roughly 1 composition in 42 to 47 at 12 branch vertices. A level is claimed only when a separate tool, run from a frozen code snapshot, does three things: it re-derives every work unit from a hashed census manifest, rebuilds every pruning table, and matches the composition counts shape by shape. The prune's arithmetic and Hall inequalities are also proved in Lean, conditional on the shape census being complete. \emph{Scope:} every non-existence result rests on exhaustive computation, not on a Lean proof.

\section{How the agents did it}
The research ran from 14 to 20 September 2026 in a ToolsEnabled agent tree. A manager agent assigned lanes to worker agents and merged their results. The workers surveyed the literature, built the shape reduction, CUDA kernels and Lean certificates, and ran the searches on two consumer GPUs, an RTX~5060~Ti desktop and an RTX~4070 laptop. Other agents reviewed the work adversarially. They wrote a referee-style report that caught overclaims, re-implemented the verifiers independently, and attacked the claim tool with forged tiers, which it rejected. The blind control above was set up to test the pipeline, and it found a witness no one had pointed it at. J.~Pinckard posed the problem, set its direction and compute limits, selected and evaluated the outputs, and is responsible for the claims. The repository holds the agents' reports, reviews, claim files and full commit history.

\vspace{0.3ex}
{\footnotesize\textbf{References.}
[1]~J.~A.~Bondy, Pancyclic graphs I, \emph{J.\ Combin.\ Theory Ser.\ B} 11 (1971) 80--84.
[2]~T.~F.~Bloom, Erd\H{o}s Problem \#1016, \url{https://www.erdosproblems.com/1016}.
[3]~S.~Griffin, Minimal pancyclicity, arXiv:1312.0274 (2013).
[4]~P.~White and Claude, Erd\H{o}s \#1016 working report (2026-07-29), \url{https://www.erdosproblemaday.com/report/1016}.
[5]~Robinfxa, JSP-000846: $m(41)=47$ (2026-09-17), \url{https://github.com/Robinfxa/JSP-000846}.
[6]~J.~C.~George, A.~Khodkar, W.~D.~Wallis, \emph{Pancyclic and Bipancyclic Graphs}, Springer (2016); not consulted here.}

\end{document}
